%\chapter*{}
\section{Introduction}
%\subsection{Preliminary remarks}

The first transport simulations came up in the late sixties when 
the problem was formulated by B.~B.~Kadomtsev and O.~P.~Pogutse \cite{KP}
and one of the first transport codes for a tokamak was created by
Yu.~N.~Dnestrovski and D.~P.~Kostomarov \cite{DK}. 
One of the authors (GVP) of this report was among the users of this
first code and is strongly influenced by the ideas implemented there. 
Further development of transport codes passed through several stages
implied by developments in plasma theory, tokamak experiments and
computing techniques. 

In spite of great progress achieved in tokamak physics during years of
research, the transport processes are still far from being really
understood. Although theoretically derived transport models are
constantly developing and their predictability improves, today's
transport modeling is to a large extent empirical and based on scaling
studies and trial-and-error approach. 
The latter assumes that a number of calculations should be accomplished
until a reasonable agreement between experimental data and modeling
results is achieved. 

The transport code used for this goal should be easily tunable and
adjustable for new tasks arising during the study.  
On the other hand, the set of equations to be solved numerically is
strongly coupled: the same physical quantity appears in many different
places in the code, so that even a minor change can invoke a long chain
of related modifications.  
This feature suggested the value of an automated code builder, which was
implemented in the first version of the Astra code at the late eighties
in Kurchatov Institute in Moscow. 

The current version has benefitted from years of experience in transport
modeling, and comprises of the majority of tools which are required for
the transport analysis of tokamak experiments. 
Nevertheless, the practice of transport modeling shows that new
requirements are continuously arising.  
Therefore the code is continuously being developed, and new features are
being added increasing its functionality. 
The ability of Astra to acquire new features, without modifying the
background organization and user interface, has made it one of
the most popular transport simulation tools in the fusion community. 
The code is opened for modifications and expansion into new areas. 
Dozens of people have already contributed their developments into
Astra.
The authors are thankful to all Astra users, and especially those who
have contributed useful suggestions and have participated in the
improvement of the code. 
Particular thanks are due to C.~M.~Roach, who read the manuscript and 
essentially contributed to its improvement. 

The currently supported version 5.3 of the Astra code runs on Sun, IBM,
DEC UNIX--workstations and on IBM PC (under Red Hat Linux). 
It uses UNIX C shell, Fortran/C compilers, and the X11 graphic libraries,
all of which are available on all UNIX systems. 
Installation requires about 20 MB of system disk space (this is shared
by all users in the same file system). 
In addition, the minimum requirement for each user is 10 MB, which can
however be considerably exceeded depending on the size of personal
libraries. 
The code is freely distributed and requires no special licenses.

The ASTRA (Automatic System for TRansport Analysis) code is more than a
transport code in the conventional sense. 
It is a flexible programming system capable of creating numerical codes
for predictive or interpretative transport modeling, for stability
analysis or for processing experimental data. 
It is also not restricted to tokamak applications exclusively.
Some features are added to versions 3.0 and onwards which allow for 1D
stellarator modeling.  
The Astra system comprises
\begin{itemize}
\item[--] an extensive library of modules describing different physical
processes and data treatment
\item[--] a supervising shell which keeps track of changes in the
libraries, processes user's requests and assembles different modules in
a required application 
\item[--] a graphic interface enabling transparent and user-friendly run
control and flexible data presentation
\item[--] an interface to an experimental database
\item[--] a help tool with built-in descriptions and a set of examples
which facilitate use of the code. 
\end{itemize}

Transport codes generated by Astra have a modular organization so that
each physical process or derived quantity is called by name which
represents an object with encapsulated implementation. 
In the majority of cases, the user deals only with the object names and
is not involved in the implementation details.  
The modular organization allows a variety of processes easily to be 
included in the code, achieving different levels of the description of 
experimental device required by the problem under consideration. 
The encapsulated implementation of objects also decreases the 
possibility of errors by ensuring that the same expression for a physical 
quantity is used in all places of the transport code.

Another significant feature of the system is that it generates
interactive codes. 
The user can observe the time evolution of plasma parameters during the
program execution, interrupt it and change data presentation and control
parameters influencing the course of the modeling. 
This enables the testing of different transport hypotheses at run time,
and so increases the efficiency of the modeling. 
The code can also run in background mode, which is useful for routine
calculations, or for time consuming cases, e.g. when many
additional heating schemes are involved. 
All output information is then saved and can be retrieved afterwards to
study the time evolution of the discharge.

The paper is organized as follows.
Section~2 provides an extended introduction to the Astra system:
it describes the general structure of the code and gives a bird's eye
view of the code. 
The main physics background is presented in Section~3, which discusses
the geometry of magnetic surfaces and related coordinate systems,
the 2D equilibrium equation and a system of 1D transport equations and
their closure (1.5D set of equations), 
appropriate boundary and initial conditions,
possible energy and particle sources and sinks including an
overview of the main schemes of auxiliary heating and current drive.
This section also contains a number of formulae which are useful for
setting transport modeling tasks. 
While providing the general background, this section is not directly
related to the Astra code, and presents a reference for the detailed
specification of transport problems. 

Astra code variables and the adopted system of units are introduced in
Section~4 (Astra Reference Guide), where the internal representation of
the system of transport equations is also formulated. 
Further discussion in this section is devoted to the generation
of user created transport models. 
Then different modules of the code which are not directly related to
transport and 
the corresponding interfaces with the transport core are discussed. 
Section~5 (Astra User's Guide) describes the Astra modeling
language used for the creation of transport models.
This section also includes the description of the Astra experimental
database format and the transfer of database information to simulation
run. 
Starting the code, run-time operation and control are also discussed
there. 

\vfill\eject
\section{Overview of the Astra code}

There are several features, which the user has to be aware of, to
understand the operation of Astra.
First of all, Astra is not a ready-to-run application, but a tool
for building customized computer codes for the solution of a
variety of transport problems in magnetically confined plasmas.
Customization of the code is the reason Astra's for extremely high
flexibility and effectiveness.
Such a programming paradigm requires two major steps:
specification of the transport problem, and performing simulations for
selected conditions.
In addition simulations may be done in two different ways: 
depending on the nature of the problem under consideration, the user may
select to run simulations in the active control mode or the background
mode. 
The description of the user interaction with Astra is shown in the
{\bf high level usage diagram}.

\begin{picture}(400,360)
\put(100,310){\begin{picture}(50,50)%
		\put(0,20){\oval(80,40)}
		\put(-16,15){{\large\bf User}}
	      \end{picture}}
\put(300,310){\begin{picture}(50,50)%
		\put(0,20){\oval(80,40)}
		\put(-18,15){{\large\bf Astra}}
	      \end{picture}}
\put(100,310){\vector(0,-1){20}}
\put(40,277){\begin{picture}(50,50)%
		\put(0,0){\fbox{Specifies transport task}}
		\put(60,-27){\line(0,1){20}}
		\put(60,-27){\vector(1,0){140}}
	      \end{picture}}
\put(240,247){\begin{picture}(50,50)%
		\put(0,0){\fbox{Builds computer code}}
		\put(70,-27){\line(0,1){20}}
		\put(70,-27){\vector(-1,0){140}}
	      \end{picture}}
\put(40,217){\begin{picture}(50,50)%
		\put(0,0){\fbox{\parbox{120pt}{Specifies machine 
				para\-meters and starts run}}}
		\put(60,-33){\line(0,1){20}}
		\put(60,-33){\line(1,-1){20}}
		\put(60,-33){\line(-1,-1){20}}
		\put(60,-73){\line(1,1){20}}
		\put(60,-73){\line(-1,1){20}}
		\put(46,-56){{\it mode}}
		\put(0,-50){{\it active}}
		\put(86,-50){{\it background}}
		\put(80,-53){\vector(1,0){120}}
		\put(40,-53){\line(-1,0){50}}
		\put(-10,-53){\line(0,-1){70}}
		\put(-10,-123){\vector(1,0){210}}
	      \end{picture}}
\put(240,160){\begin{picture}(50,50)%
		\put(0,0){\fbox{\parbox{110pt}{Performs simulations 
					and saves results}}}
		\put(70,-33){\line(0,1){20}}
		\put(70,-33){\vector(-1,0){170}}
		\put(-175,-37){\fbox{Views results}}
	      \end{picture}}
\put(240,75){\begin{picture}(50,50)%
		\put(0,0){\fbox{\parbox{112pt}{Performs simulations\\ 
					and displays evolution\\ of parameters}}}
		\put(0,-13){\vector(-1,0){77}}
		\put(-77,-7){\vector(1,0){77}}
	      \end{picture}}
\put(40,60){\begin{picture}(50,50)%
		\put(5,0){\fbox{\parbox{110pt}{Controls simulations\\ 
				at run time}}}
		\put(60,-33){\line(0,1){20}}
		\put(60,-33){\vector(1,0){140}}
		\put(200,-35){\fbox{\parbox{110pt}{Saves and prints data
				if requested}}}
	      \end{picture}}
\end{picture}

%            ------------                   -------------
%            |   User   |                   |   Astra   |
%            ------------                   -------------
%
%           Specifies
%           transport task
%                   |
%                   ----------------------> Builds
%                                           computer code
%                                               |
%           Specifies machine <-----------------
%           parameters and
%           starts run
%                   |
%                   |
%       active     / \  background          Performs
%        --------- mode ------------------> simulations and
%        |         \ /                      saves results
%        |                                          |
%        |                                          |
%        |       Views results    <------------------
%        |
%        |
%        ---------------------------------> Performs
%                                           simulations and
%            Controls simulations <------   displays evolution
%            at run time           ------>  of parameters
%                        |
%                        |
%                        ---------------->  SAVES and print data
%                                           if requested
%

Let us now consider each step of operation with Astra in more detail.
To set the transport task the user basically has to specify how to
perform the simulations and what to display.
Some adjustment of these settings may be done through run time control,
but the overall simulation tasks of the generated code are determined
during the initial transport model specification.
The following features of Astra are used to specify the model description:
\\ \hspace*{5 mm} -- transport modeling language
\\ \hspace*{5 mm} -- standard expressions for physical quantities
\\ \hspace*{5 mm} -- modules representing physical processes or
features of the experimental device. \\
The modeling language is described in Section~5.2, standard expressions
are listed in Sections~4.7 and 4.8, basic modules are discussed in
Section~4.9. 

To convert the model description into an executable code, Astra
interprets the model description (source prototype) and compiles the
generated source code. 
The sequence of operations and Astra objects involved in the whole
process of creating the code and obtaining simulations results are shown
in the {\bf flow diagram}.\label{flow}

\begin{picture}(400,60)
%    Language            Expressions             Modules
%        \                    |                    /
%          \                  |                  /
%            \                |                /
%             ---------------------------------
%             | User specifies transport task |
%             ---------------------------------
\put(-20,10){\hbox{\sc Step 1}}
\put(20,10){$\lb\{\mbox{\begin{picture}(50,40)
		\put(0,25){\fbox{Language\strut}}
		\put(90,25){\fbox{Expressions\strut}}
		\put(200,25){\fbox{Modules\strut}}
		\put(25,15){\vector(3,-1){64}}
		\put(125,15){\vector(0,-1){22}}
		\put(225,15){\vector(-3,-1){64}}
		\put(70,-25){\fbox{{\bf User specifies task\strut}}}
	    \end{picture}}\rb.$}
%                             |
%                             |
\put(159,-30){\vector(0,-1){22}}
\end{picture}
%\begin{picture}(400,20)
%\put(30,170){\begin{picture}(50,50)%
%		
%	     \end{picture}}
%\end{picture}

\begin{picture}(400,200)
%                     Model description     <-------- Astra language
%                             |
% Source prototype            |                Expression lib
%           \                 |                 /
%             \               |               /
%               \             |             /
%             --------------------------------
%             |  Interpreter builds source   | ------> model report
%             --------------------------------
\put(-20,90){\hbox{\sc Step 2}}
\put(20,90){$\lb\{{\begin{picture}(50,55)
		\put(0,10){\fbox{Source prototype\strut}}
		\put(80,40){\fbox{Model description\strut}}
		\put(310,40){\fbox{\fbox{Astra language\strut}}}
		\put(300,43){\vector(-1,0){110}}
		\put(160,10){\fbox{Expression lib\strut}}
		\put(40,0){\vector(2,-1){44}}
		\put(125,30){\vector(0,-1){52}}
		\put(210,0){\vector(-2,-1){44}}
	       \put(40,-40){\fbox{{\bf Interpreter builds source\strut}}}
		\put(230,-40){\fbox{Model report\strut}}
		\put(200,-37){\vector(1,0){25}}
	     \end{picture}}\rb.$}
%                             |
%                             |
\put(159,35){\vector(0,-1){22}}
\end{picture}

\begin{picture}(400,100)
%                         Source code       <-------- FORTRAN & C
%                             |
% Build-in modules            |               Plug-in module lib
%           \                 |                   /
%             \               |                 /
%               \             |               /
%               -------------------------------
%               |  Compiler builds executable |
%               -------------------------------
\put(-20,55){\hbox{\sc Step 3}}
\put(20,55){$\lb\{{\begin{picture}(50,55)
		\put(0,10){\fbox{Built-in modules\strut}}
		\put(90,40){\fbox{Source code\strut}}
		\put(160,10){\fbox{Plug-in module lib\strut}}
		\put(40,0){\vector(2,-1){44}}
		\put(125,30){\vector(0,-1){52}}
		\put(210,0){\vector(-2,-1){44}}
		\put(315,40){\fbox{\fbox{Fortran \& C\strut}}}
		\put(305,43){\vector(-1,0){140}}
	      \put(50,-40){\fbox{{\bf Compiler builds executable\strut}}}
	     \end{picture}}\rb.$}
%                             |
%                             |
\put(159,5){\vector(0,-1){22}}
\end{picture}

\begin{picture}(400,200)
%                             |
%                             |
%                       Executable code     <-------- binary
%                             |
% Experimental database       |               Run-time control
%           \                 |                   /
%             \               |                 /
%               \             |               /
%               -------------------------------
%               | System runs simulation task |
%               -------------------------------
%                             |
%                             |
%                         results
\put(-20,105){\hbox{\sc Step 4}}
\put(20,105){$\lb\{{\begin{picture}(50,85)%
		\put(125,107){\vector(0,-1){22}}
		\put(85,65){\fbox{Executable code\strut}}
		\put(330,65){\fbox{\fbox{Binary\strut}}}
		\put(325,68){\vector(-1,0){140}}
		\put(125,55){\vector(0,-1){57}}
		\put(0,30){\fbox{Experimental data\strut}}
		\put(150,30){\fbox{Run-time control lib\strut}}
		\put(40,20){\vector(2,-1){44}}
		\put(205,20){\vector(-2,-1){44}}
		\put(40,-20){\fbox{{\bf System runs simulation task\strut}}}
		\put(125,-30){\vector(0,-1){22}}
		\put(105,-70){\fbox{Results\strut}}
		\put(320,-70){\fbox{\fbox{\parbox{20mm}
			{Graphics\\ ASCII\\ PostScript}}}}
		\put(315,-70){\vector(-1,0){155}}
	   \end{picture}}\rb.$}
\end{picture}
%    Language            Expressions             Modules
%        \                    |                    /
%          \                  |                  /
%            \                |                /
%             ---------------------------------
%             | User specifies transport task |
%             ---------------------------------
%                             |
%                             |
%
%                     Model description     <-------- Astra language
%                             |
% Source prototype            |                Expression lib
%           \                 |                 /
%             \               |               /
%               \             |             /
%             --------------------------------
%             |  Interpreter builds source   | ------> model report
%             --------------------------------
%                             |
%                             |
%                         Source code       <-------- FORTRAN & C
%                             |
% Build-in modules            |               Plug-in module lib
%           \                 |                   /
%             \               |                 /
%               \             |               /
%               -------------------------------
%               |  Compiler builds executable |
%               -------------------------------
%                             |
%                             |
%                       Executable code     <-------- binary
%                             |
% Experimental database       |               Run-time control
%           \                 |                   /
%             \               |                 /
%               \             |               /
%               -------------------------------
%               | System runs simulation task |
%               -------------------------------
%                             |
%                             |
%                         results
%
\\[-15pt]
Specifying the simulation task is still a programming job, but it is
simplified greatly by using Astra's high-level transport problem
oriented language. 
To add a specific physical process to the model description the user has
only to call it by name in the proper context. 
Analyzing these requests the Astra interpreter creates a source code
where the expressions corresponding to the process names are used in the
way specified by the context. 
These expressions are stored in the expression library and may be viewed
or modified by the user. 
Similar things happen with modules describing auxiliary properties of
experimental device such as heating or current-drive systems. 
Even if these modules involve a lot of functionality and controls, the
user may easily add them to or remove them from his model description at
his own discretion depending on the specifics of the transport task. 
Due to this add/remove option these modules are called plug-in. 
In addition, transport codes generated by Astra contains a number of
built-in modules such as the time evolution driver, a graphical
representation of output, a supervising control shell, and other features 
of the transport task which are always included in the simulation model.
This part of the code provides the framework for the operation of the
customized part and is hidden from the ordinary user. 
Modification of this part is rarely required, and an increased Astra
functionality is available through the use of expression and plug-in
module libraries. 

A few words have to be said about executable transport codes created by
Astra.
Though executable code is customized for each specific task it always
has the following set of standard features:\\
\hspace*{5mm} -- it may be run either in the active or background
		modes, depending on the type of \linebreak
		\hspace*{9mm} simulations performed by the user\\ 
\hspace*{5mm} -- in the active mode, code enables user specified 
			run time control\\ 
\hspace*{5mm} -- the code interfaces with the experimental database 
			which includes discharges from %\linebreak
			\hspace*{9mm} many machines\\ 
\hspace*{5mm} -- it has a modular organization.\\
The first three features are represented in the high level usage diagram.
They will be discussed in more detail in Section~5.
Modular organization is supported by the method of the code creation.
Each block necessary for simulations is connected to other blocks
through the Astra environment providing uniform control and exchange
capability. 
Plug-in modules are included and connected through the same mechanism.
Modular organization allows easy modification of the code and provides
an efficient tool for building transport codes to any level of detail. 
The generated transport code has a transparent image in model
description (prototype) and summarized in the model report, 
which in a compact form and without any omissions 
represents the whole complexity of the executable transport code.

The structure of the Astra code is shown in the {\bf module diagram}.
\\
\begin{picture}(400,120)
   \put(100,90){\fbox{\fbox{Astra environment and internal
			data storage\strut}}}
   \put(130,75){\vector(-3,-1){95}}
   \put(140,75){\vector(-1,-1){32}}
   \put(185,75){\vector(0,-1){32}}
   \put(260,75){\vector(0,-1){32}}
   \put(310,75){\vector(1,-1){32}}
   \put(330,75){\vector(3,-1){95}}
   \put(0,20){\fbox{\parbox{60pt}{Interface to\\ database\strut}}}
   \put(85,20){\fbox{\parbox{45pt}{Built-in\\ modules\strut}}}
   \put(155,20){\fbox{\parbox{53pt}{Transport\\ equations\strut}}}
   \put(235,20){\fbox{\parbox{45pt}{Plug-in\\ modules\strut}}}
   \put(310,20){\fbox{\parbox{65pt}{Data display and control\strut}}}
   \put(405,20){\fbox{\parbox{45pt}{Output\\ interface\strut}}}
\end{picture}
%
%           -----------------------------------------------
%           | Astra environment and internal data storage |
%           -----------------------------------------------
%          /      |          |          |         |        \
%        /        |          |          |         |          \
%      /          |          |          |         |            \
%Interface     Built-in   Transport  Plug-in   Data display   Output
%to database   modules    equations  modules   and control    storage
%

The other feature of Astra, which has to be discussed, is that Astra is a 
multi-user distributed system. 
Generally, there are the code kernel and several working directories 
for each user. 
The kernel of the code is shared by all users through a set of symbolic 
links to the kernel directory. 
This enables user's access to libraries of expressions, plug-in 
modules and object modules.
On the other hand, all core parts of the code, such as the model
interpreter, supervising shell, code building tools and built-in modules
are present in the kernel only and consequently cannot be modified by an
unauthorized user. 
This shared configuration saves disk space and more importantly keeps
all users with the most recently updated version of the kernel while
only one version of the kernel needs to be maintained. 

The user owns all variable parts of the code such as model descriptions, 
executable codes, experimental libraries and the results of modeling.
A trickier issue is the ownership of the libraries of standard
expressions and physical processes.
Basically they belong to the kernel directory and a user has read-only
access to the sources files.
This allows viewing and linking standard modules in a user created code
but forbids any changes in these modules.
The latter is however not really restrictive since every user
can have his own library and use it together with the common one.
Users may also create new modules or expressions for their own
particular needs. 
In particular, a user can copy a common module to his own directory 
with a new name, then include it in his personal library 
thus acquiring all rights for subsequent modifications. 

Physical locations of different blocks of the Astra system 
and scheme of their interaction are shown in
the {\bf component diagram}.\\%[-10pt]
\begin{picture}(400,380)
\put(0,320){\begin{picture}(50,50)
		\put(0,40){{\parbox{120pt}{\sc Kernel directory}}}
		\put(200,40){{\parbox{320pt}{\sc
			User's Astra working directory 
			}}}
		\put(-5,25){$\overbrace{\mbox{\hspace*{110pt}}}$}
		\put(140,25){$\overbrace{\mbox{\hspace*{320pt}}}$}
		\put(0,-10){\fbox{\parbox{100pt}{\begin{center}\bf
			Common library 
			\end{center}}}}
		\put(150,-10){\fbox{\parbox{100pt}{\begin{center}\bf
			User library 
			\end{center}}}}
		\put(50,-35){\vector(0,-1){35}}
		\put(200,-35){\vector(-4,-1){140}}
	    \end{picture}}
\put(0,220){\begin{picture}(50,50)%
		\put(0,0){\fbox{\parbox{100pt}{\begin{center}\bf
			Code builder 
			\end{center}}}}
		\put(150,0){\fbox{\parbox{100pt}{\begin{center}\bf
			Transport code\\ prototype
			\end{center}}}}
		\put(310,0){\fbox{\parbox{150pt}{\begin{center}\bf
		Device parameters\\ Experimental data\\ Initial conditions
			\end{center}}}}
		\put(50,-25){\line(0,-1){70}}
		\put(50,-95){\vector(1,0){85}}
		\put(145,3){\vector(-1,0){32}}
		\put(390,-34){\line(0,-1){13}}
		\put(390,-47){\line(-1,0){190}}
		\put(200,-47){\vector(0,-1){13}}
	     \end{picture}}
\put(0,120){\begin{picture}(50,50)%
%		\put(200,5){\oval(130,60)}
		\put(140,0){\fbox{\parbox{120pt}{\begin{center}\bf
			Executable code for\\ transport data analysis 
			\end{center}}}}
		\put(280,20){\it Graphic}
		\put(280,10){\it ~output}
		\put(280,-10){\it Run--time}
		\put(280,-20){\it ~control}
		\put(385,3){\oval(90,60)}
		\put(325,0){{\parbox{120pt}{\begin{center}\bf
			Interactive\\ control panel 
			\end{center}}}}
		\put(350,-40){\line(1,0){70}}
		\put(360,-30){\line(1,0){50}}
%		\put(340,-60){\line(1,0){80}}
		\put(350,-40){\line(1,1){10}}
		\put(420,-40){\line(-1,1){10}}
		\put(200,-35){\vector(0,-1){20}}
		\put(270,5){\vector( 1,0){65}}
		\put(335,1){\vector(-1,0){65}}
	     \end{picture}}
\put(0,33){\begin{picture}(50,50)%
		\put(0,0){\fbox{\parbox{100pt}{\begin{center}\bf
			Post--viewer 
			\end{center}}}}
		\put(150,0){\fbox{\parbox{100pt}{\begin{center}\bf
			Stored output data 
			\end{center}}}}
		\put(145,3){\vector(-1,0){32}}
	     \end{picture}}
\end{picture}
%	here: (width,height)	1pt = 0.35mm
%\begin{picture}(150,450)
% Environment {picture} allows to put few figures in a line,
% 				otherwise they are put in column
% epsfysize = picture_height {\epsfbox[ix1 iy1 PSwidth PSheight]{...}}
% epsfysize is height used by a picture
% (ix1,iy1) - position of image relative to the allocated box (picture)
% 	      the reference points are lower left corners of both
% (PSwidth,PSheight) region of PS to be plotted
% Picture is then scaled so that (iy1+PSheight) fits epsfysize
%\epsfysize 15cm {\epsfbox[0 150 540 720]{astra1.PS}}
%\epsfysize 10cm {\epsfbox[0 0 530 715]{astra1.PS}}
%\begin{picture}(150,450)
%\epsfysize 15cm {\epsfbox[0 150 540 720]{astra1.PS}}
%\end{picture}
\\[-10pt]
At the code installation, the user's Astra directory structure is
created. 
We denote the root directory of this structure as \tw{AWD} (Astra
Working Directory). 
This directory is a current working directory during the code execution. 
In the code, relative addresses only are used, therefore, \tw{AWD}
can have arbitrary location with respect to user's home directory.
Moreover, the user can simultaneously handle several versions of the
Astra code, each linked to different kernels.
However, if another name is not given explicitly then \tw{AWD} is
$\tt\symbol{36}HOME/astra$. 
Specific components are stored in subdirectories which are given in
Table 2.1:
\medskip\\
\indent
{\bf Table 2.1: Astra working directory structure}
\nopagebreak
\\[10pt]
\nopagebreak
\noindent%p{1.5cm}
\begin{tabular}{| l | l | l | l |}				\hline
\bf Directory&\bf Contents	
			&\bf File owner&\bf\hspace*{-1.5pt}User access\\ \hline
\tw{AWD/equ/}	& Transport code prototypes and \tw{*.log} files
			& User & Allowed 	\\
\tw{AWD/exp/}	& Device parameters, experimental data base 
			& User & Allowed  	\\
\tw{AWD/udb/}	& Input and output data in the form of u-files 
			& User & Allowed  	\\
\tw{AWD/dat/}	& Computation results (\tw{ASCII}) 
			& User & Allowed  	\\ \hline
\tw{AWD/fml/}	& Simple expressions (formulae)
			& Mixed & Allowed 	\\
\tw{AWD/fnc/}	& Extended expressions (\tw{FORTRAN} functions)
			& Mixed & Allowed 	\\
\tw{AWD/sbr/}	& Plug-in modules (\tw{FORTRAN} subroutines)
			& Mixed & Allowed 	\\ \hline
\tw{AWD/tmp/}	& Included intermediate \tw{FORTRAN} source files 
			& User & Restricted 	\\
\tw{AWD/.res/}	& Computation results (post-viewer binary files) 
			& User & Restricted 	\\
\tw{AWD/.tsk/}	& Executable transport codes 
			& User & Restricted 	\\ \hline
\tw{AWD/for/}	& Service and built-in modules 
			& System & Not allowed 	\\
\tw{AWD/.lbr/}	& Common libraries and executable modules 
			& System & Not allowed	\\
\tw{AWD/.exe/}	& Core and built-in modules 
			& Mixed & Not allowed	\\ \hline
\end{tabular}
%\\[10pt]
\vskip 10pt

In most practical work, the user deals mainly with the first 4 directories 
of this table. 
All these directories include user owned files in ASCII format only.
The directory \tw{AWD/equ/} contains transport code prototypes 
(in what follows ``models'') and related files. 
The syntax of model files is discussed in Section~5.2.
The directories \tw{AWD/exp/} and \tw{AWD/udb/} comprise
of the experimental data base for interpretative modeling and
initial data for predictive simulations.
Astra format for input files is described in Section~5.4.
Computation output files, either in ASCII or in Post Script format, are
stored in \tw{AWD/dat/}. 
Output data files have self-evident form and do not require explanations. 

More advanced users may also add files to directories
\tw{AWD/fml/}, \tw{AWD/fnc/}, \tw{AWD/sbr/}.
These directories can include both shared (soft links) and personal files.
Every new file appearing in these directories is automatically
compiled and added to a user library of object modules, so that
it can then be used in the transport code.

The last six directories in the table are used by the Astra system
but the unauthorized user has no access to change files in these
directories.
This part of the Astra system and its kernel are not described 
in the current report.

\vfill\eject
